\documentclass{article}
\usepackage[utf8]{inputenc}
\usepackage{changepage}% http://ctan.org/pkg/changepage
\usepackage{float}
\usepackage{fancyhdr}
\usepackage{lastpage}
\usepackage{graphicx}
\usepackage{ragged2e}
\usepackage{scrextend}
\usepackage{lastpage}
\usepackage{geometry}
 \geometry{
 left = 1.5in,
 right = 1.5in,
 top = 1in
 }
\pagestyle{fancy}
\renewcommand{\headrulewidth}{0pt}
\rhead{\includegraphics[height= 1cm]{ento1.png}\vspace{1em}}
\chead{\includegraphics[width = 5cm]{agritech1.png}}
\lhead{\includegraphics[width=3cm]{acet1.png}}
 
\fancyfoot{}
\setlength\headheight{40.8pt}
\RequirePackage[colorlinks=true, allcolors=blue]{hyperref}
\usepackage{titlesec}
\titleformat{\section}{\normalfont\fontsize{10}{15}\bfseries}{\thesection}{0em}{}
\titlespacing{\section}{}{1.24em}{0.24em}
 
\usepackage[default]{lato}
\usepackage[T1]{fontenc}

 
\cfoot{Page \thepage \hspace{1pt} of \pageref{LastPage}}

\begin{document}

\begin{addmargin}[3.8in]{}
Willett, Filgueiras, \newline
Nyrop and Nault \newline
----------------------------- \newline
Denis S. Willett\newline
Department of Entomology \newline
Cornell AgriTech  \newline
Cornell University \newline
15 Castle Creek Drive \newline
Geneva, NY 14456 \newline
deniswillett@cornell.edu \newline
\end{addmargin}
\setlength{\parindent}{0cm}

December 24, 2019

\vspace{1.24em}

Pest Management Science

\vspace{1.24em}

Dear Dr. Torto,

\vspace{0.48em}
\setlength{\parindent}{1.24cm}

Thank you for your consideration of \textit{“Attract and kill: Spinosad containing spheres to control onion maggot (\textit{Delia antiqua})”} for publication in Pest Management Science.  We have revised the manuscript in accordance with reviewer recommendations.  A complete and detailed account of the revisions is included below.  The reviewer comments were very helpful, much appreciated, and have resulted in a better manuscript.  In addition to submitting a revised version of the manuscript incorporating the latest changes, we are also submitting production quality images.  
  

\vspace{2em}

Much appreciated, \newline

\vspace{1em}

Denis S. Willett
\vspace{0.48em}

Camila C. Filgueiras
\vspace{0.48em}

Jan P. Nyrop
\vspace{0.48em}

Brian A. Nault
\vspace{0.48em} 


\vspace{1.4em}

\hline

\setlength\parindent{0pt}

\subsection*{Referee 1}

The authors did a great job writing the manuscript.  The introduction, methods, and results were easy to understand.    

\begin{quote}
    \textit{Thank you for reading the manuscript and providing constructive feedback.  We appreciate the suggestions and have endeavored to include them in the revised version.  }
\end{quote}

On lines 54 and 55 the authors state field trials were conducted from 2006 to 2008, but they did NOT describe if all experiments were repeated in 2006, 2007, and 2008.  The authors need to describe which experiments were conducted each year, how many times each experiment was replicated, and if the experiments were replicated in space and time.  

\begin{quote}
    \textit{Thanks for the opportunity to clarify.  Details on the timing of trials are now included in at the beginning of each sub-section in the Material and Methods section on lines ....  In short, the trials were replicated in space and time.  }
\end{quote}


The authors should add a few sentences describing the field sites (lines 54-56) cropping history (e.g. Where onions grown in previous years?) and the reason why the fields were chosen (e.g. did the fields have a previous history of onion maggot damage?).  If experiments were not replicated in space or time it should be noted in the Material and Methods. 

\begin{quote}
    \textit{Thanks for mentioning this opportunity to elaborate.  Details on cropping history and reasoning behind field site choice have now been included on lines ....  Information on how experiments were replicated in space and time are also included (see above).  }
\end{quote}

In the discussion section the authors made no mention of using spinosad seed treatment as a management option in conventional and organic onion fields.  Why?? Spinosad (Regard) seed treatment is currently the primary management method for onion maggot in organic onions, and spinosad is an active ingredient in many conventional seed treatments.   Spinosad is an active ingredient in the FI500 seed treatment, and Regard alone often provides a similar level of protection compared to FI500.  Since FI500 resulted in lower maggot damage compared to all "attract and kill" tactics, why would someone recommend attract and kill with spinosad over spinosad seed treatment?   I strongly recommend the authors discuss the advantages, disadvantages, and compatibility of using spinosad seed treatment and spinosad attract & kill.    

\begin{quote}
    \textit{This is an excellent point.  We've brought this up and elaborated more on exactly this point in the discussion as well as addressing the advantages and disadvantages of these strategies.  Thanks for the suggestion.  Additions can be found on lines ....
    }
\end{quote}

The authors should discuss the implications of onion maggot developing resistance to spinosad using the attract and kill tactic.  e.g. Will the attract and kill tactic with spinosad lead to a higher likelihood of onion maggot developing resistance to spinosad?  This tactic targets egg laying adults, it only provides partial control, and it exposes multiple generations per year to spinosad?   These factors seem to favor resistance. This question is especially relevant as the most popular registered insecticide seed treatments contain spinosad and there are few alternatives in the marketplace.     

\begin{quote}
    \textit{Also an excellent suggestion.  This is particularly important in systems where options are rapidly disappearing.  We've addressed this in more detail in the discussion on lines ...
    }
\end{quote}

\subsection*{Referee 2}

The methods and results need some significant improvements. I found that methods lacking in details and the analysis not well described. I also noted in the annotated file that the quality of the writing in the intro and discussion is markedly different and will need some work to bring them to the same level.

\begin{quote}
    \textit{Thank you for your comments regarding the manuscript.  We have added extended details to the methods section describing both the field work and analysis.  Additionally, we have adjusted language in the introduction and discussion.  We appreciate your feedback in helping improve the manuscript.  }
\end{quote}

Line 60: Please check and correct.  

\begin{quote}
    \textit{Thanks for the opportunity to verify this.  It is correct.  'Abut' (spelled exactly that way) means to share a common border. }
\end{quote}

Line 63: Not sure why its assumed that color and rate are related, explain.  

\begin{quote}
    \textit{Interesting question.  There is no assumption that color and rate are related.  The fact that white spheres are attractive to \textit{D. antiqua} is based on previous research in this system.  }
\end{quote}

Line 65: Does this sentence belong to the next paragraph? move it there.

\begin{quote}
    \textit{Excellent suggestion.  This sentence does seem to be hanging there.  Rather than include it in the subsequent paragraph, we have opted to remove it alltogether.  }
\end{quote}

Line 66: be more specific.  

\begin{quote}
    \textit{Thanks for the opportunity to elaborate.  We have included more details as to the timing and location of the trials throughout the methods section, specifically lines .... }
\end{quote}


Line 68: climate or weather?

\begin{quote}
    \textit{Interesting question.  We want to use the weather in these trials as a proxy for general exposure to northern climates and have revised the sentence in accordance with your suggestion.  }
\end{quote}


Line 68-69: This was already state above, seems redundant.

\begin{quote}
    \textit{Thanks for mentioning this.  We have removed the redundant statement.  }
\end{quote}

Line 71: More details are needed: how many reps/field, how many fields, how much distance between spheres, etc. Where in the field relative to the crop were spheres? What were they placed on and how high?

\begin{quote}
    \textit{Thanks for the opportunity to elaborate.  Details about reps, fields, distances, location, and height have been included on lines ....}
\end{quote}

Line 77: Which years was this experiment done? Throughout this section pls change to using SI units.

\begin{quote}
    \textit{Thanks for the suggestions.  Information on the timing of the experiment has been added to line.... Units have also been changed to SI.  Details have also been added regarding replications and field sites (lines ...).  }
\end{quote}

Line 91: was this based on a preliminary experiment?

\begin{quote}
    \textit{This was indeed based on previous work establishing the attractiveness of the spheres (Willett et al 2019). }
\end{quote}

Line 99: novel? How was it different from previously used ones?

\begin{quote}
    \textit{New but not novel.  It differed from previous formulations with the addition of the two extra components included below.  We have elaborated on the reasons for developing this on lines.   }
\end{quote}

Line 100: why were these compounds chosen?

\begin{quote}
    \textit{Excellent question.  The thought was that these compounds would help sterilize the paraffin mixture and enhance performance of the spheres over the long term.  }
\end{quote}

Line 104: was this true for previous experiments as well? If yes, please mention above.

\begin{quote}
    \textit{This was indeed true for previous experiments as well.  Thanks for the suggestion; these details have been included in the sections above as well.  }
\end{quote}

Line 108: which years?

\begin{quote}
    \textit{Great question.  This information has been included and elaborated upon in all methods sections on lines...}
\end{quote}

Line 110: pls explain how it was analyzed in addition to how it was displayed.

\begin{quote}
    \textit{Thanks for your inquiry.  It was analyzed using a smoothing function consisting of a generalized additive model with cubic splines.  }
\end{quote}

Line 112: aren't these the same thing?

\begin{quote}
    \textit{Thanks for your question.  Yes and No.  While linear models and analysis of variance use the same underlying model structure, they have two different objectives and answer two different questions.  Linear models represent the relationship between independent and dependent variables.  Analysis of variance compares inter and intra-group variation to examine differences.  In our case, both approaches were used because we are interested both in modeling relationships and differences between groups.  }
\end{quote}

Line 113: I'd like to get more details about which factors you have used, fixed and random.  

\begin{quote}
    \textit{Thanks for suggesting this inclusion.  We've listed the fixed and random factors used in the analysis on lines ...}
\end{quote}

Line 115: AIC?, BIC? How did you use these to inform model selection? Considerations: third time this word is used in this sentence.  

\begin{quote}
    \textit{Good question.  Both were used to avoid overfitting.  As both of these criteria penalize likelihood estimates with increasing factors, these were used in consideration with other factors listed to determine best fit models.  Redundant uses of the word 'considerations' have been removed from the sentence.  
    }
\end{quote}

Line 116:  I suggest rewording this sentence. I'm assuming that you are talking about means separation here?

\begin{quote}
    \textit{Exactly!  We've reworded this sentence to clarify.  Thanks for the suggestion.  }
\end{quote}

Line 116: Why use both, why not just one or the other? 

\begin{quote}
    \textit{Tukey HSD test generally is used to evaluate all pairwise comparisons.  Dunnett's test evaluates treatments versus controls.  In some cases we are only interested in comparisons versus controls (more powerful statistically as the correction for the family wise error rate is lower for less comparisons) while in others we are interested in more detailed comparisons.  }
\end{quote}

Line 117: Be more explicit about which model goes with what data. I'd like to see a list of dependent and independent variables and fixed and random effects.I suggest removing such a blanket statement and being more specific as I mentioned above.

\begin{quote}
    \textit{Thanks for the suggestion.  We've included more detail about the models and factors used for each datasets on lines ...  In addition, the exact code to run the analysis with annotations and extensive documentation is publicly available on GitHub.  }
\end{quote}

Line 124: need to mention these where you talk about how you built and selected models. These are the tools you used to do that.

\begin{quote}
    \textit{Thanks for the suggestion.  Given that these tools were used throughout the analysis, we thought it best to consolidate them here rather than repeatedly mentioning them throughout the text.  }
\end{quote}

Line 128: I'm not following why showing this data is relevant. You also have not done analysis on it. Suggest incorporating it into the rest of the MS effectively.

\begin{quote}
    \textit{Thanks for the opportunity to elaborate.  These data are important as they document three generations in upstate New York.  This is crucial because it highlights the importance of targeting the first generation to prevent further damage throughout the season.  While the these data are often mentioned, this figure is important to practitioners in New York state for reference and management decisions.  Generalized additive models with cubic splines were used to highlight trends in data for management professionals.  }
\end{quote}

Line 129: What period was this? Be more specific

\begin{quote}
    \textit{Excellent opportunity to clarify.  Thanks.  We've added details to lines ...}
\end{quote}

Line 131: delete "roughly"

\begin{quote}
    \textit{Excellent suggestion.  'Roughly' has been deleted.  }
\end{quote}

Line 133: What do you mean? Significantly different? Different from what?

\begin{quote}
    \textit{Great opportunity to clarify.  This has been explained on lines...}
\end{quote}

Line 143: Was this peak significantly different from the other times?

\begin{quote}
    \textit{Interesting question.  While we did not specifically test for this, peaks were either markedly reduced or noticeably absent for subsequent times in Figure 3.  }
\end{quote}

Line 154: This is compared to what?

\begin{quote}
    \textit{This comparison was made to spinosad spheres alone - specifics have been addded to lines ...}
\end{quote}

Line 155: I have a hard time with the letters above bars on this figure. I didn't see the darker colored 'A's. Usually 'A' is used for the highest numbers.

\begin{quote}
    \textit{Thanks for this suggestion.  We've updated the colors to make this darker.  The lettering scheme was chosen to read alphabetically from left to right.  }
\end{quote}

Line 155: Was this significant?

\begin{quote}
    \textit{Great question.  It was indeed significant.  Details have now been included in lines ...}
\end{quote}

Line 157: be more specific about what treatments.  

\begin{quote}
    \textit{In this case treatment was with the attract and kill strategies.  Details have been added to lines ...}
\end{quote}

Line 159: are there any F-values or model results to go along with this?

\begin{quote}
    \textit{Thanks for the opportunity to add more detail.  Additional model statistics have been included on lines ...}
\end{quote}

\textbf{
Line 162: The discussion and introduction seems to have been written by different people. The language used in the discussion is much more "fuzzy" than in the intro. I suggest re-wording much of the discussion to make it less wordy and repetitive and place the findings into a larger context.
}

Line 167: was this personal observation or speculation?

\begin{quote}
    \textit{Excellent question.  This was a personal observation. Details to that effect have been included on lines ...}
\end{quote}

Line 169: Suggest simplifying this paragraph, its a long and convoluted way of saying that dose didn't matter.

\begin{quote}
    \textit{\textbf{EDIT HERE}}
\end{quote}

Line 174: Its not clear what this sentence is about, please re-word to clarify.

\begin{quote}
    \textit{Thanks for the suggestion.  We've reworded the sentence to make it clear that we are estimating numbers of flies killed to place it in the larger context of population management.  }
\end{quote}

Line 176: caught how and where?

\begin{quote}
    \textit{Excellent opportunity to clarify.  This was clarified later in the paragraph, but merits inclusion earlier.  The adults were caught on sticky traps.  The phrase has been updated on lines...}
\end{quote}

Line 180: how does this compare to the many flies that are in a typical onion field? I'm not understanding the value of your analysis. Please clarify and related this info to a larger context.

\begin{quote}
    \textit{Thanks for your question.  This directly relates to your inquiry about Figure 1 previously.  In this paragraph, we place mortality due to the attract and kill strategy in the context of population dynamics in a typical onion field through comparison of data from the population monitoring trial discussed previously.  }
\end{quote}

Line 187: Or they were attracted but didn't contact the spheres?

\begin{quote}
    \textit{Thanks for your inquiry.  In this case, the discussion is regarding the use of sticky cards.  The previous part of this phrase explains exactly that the kill part of the strategy with sticky cards was ineffective.  }
\end{quote}

Line 188: a bit of discussion on why you suspected that these would change attractiveness would be great.

\begin{quote}
    \textit{Excellent suggestion.  We've included more details about our rationale on lines ...}
\end{quote}

Figure 3: suggest removing the background.  

\begin{quote}
    \textit{Thanks for the helpful suggestion.  We've removed the background to enhance visibility.  }
\end{quote}

Line 196: Feels repetitive, the previous sentence covered this.

\begin{quote}
    \textit{Thanks for the suggestion.  These are actually two different cases.  Resistant populations and organic growers may elect to use this strategy, but have the option to use other methods.  There are some cases where conventional pesticides to combat onion maggot have been banned or are in the process of becoming banned which would render them unavailable as a control option.  }
\end{quote}

Figure 4: light grey is difficult to see, suggest using bars instead of dots and lines.

\begin{quote}
    \textit{Excellent suggestion.  We've adjusted the color scheme in the figures to enhance readability.  }
\end{quote}


\end{document}
